%=====================================================================
\chapter{Planning an E-Commerce Framework}\label{ch:planning}
%=====================================================================
\locator{1}{Part I --- E-Commerce and Its Models}

\begin{outcomes}
By the end of this chapter you will be able to:
\begin{tight}
  \item \textbf{Distinguish} the four ways a shop can be obtained, and
        \textbf{state} what each costs in money, time and control.
  \item \textbf{Write} requirements for a shop in a form that discriminates
        between those four rather than describing all of them equally.
  \item \textbf{Identify} the small number of requirements that actually decide
        the choice, and \textbf{explain} why most requirements do not.
  \item \textbf{Compute} a total cost of ownership over three years, including
        the costs that do not appear on a price list.
  \item \textbf{Defend} a build-buy-or-extend decision against the obvious
        objection to it.
\end{tight}
\end{outcomes}

\begin{prereq}
\chref{models} for the business and revenue model, and its laboratory --- this
chapter turns that one-page brief into a decision. \chref{auctions} if your
venture is a marketplace rather than a shop.
\end{prereq}

\begin{keyterms}
requirement $\bullet$ functional requirement $\bullet$ non-functional
requirement $\bullet$ SaaS $\bullet$ hosted platform $\bullet$ self-hosted
$\bullet$ open source $\bullet$ bespoke development $\bullet$ total cost of
ownership $\bullet$ vendor lock-in $\bullet$ extensibility $\bullet$ data
portability $\bullet$ decision matrix $\bullet$ minimum viable product
\end{keyterms}

\section{The Question This Chapter Answers}

You have a business model, a revenue model and some idea of what you are
selling. There is now exactly one decision to make before any of the rest of
this book applies, and it is this: \emph{where does the software come from?}

There are four answers, and the whole of Parts II to V is a consequence of
which one you pick.

\dsfig{%
\begin{tikzpicture}[
  opt/.style={rounded corners=3pt, draw=#1, line width=1.1pt, fill=#1!8,
              text=#1!60!black, font=\bfseries\scriptsize, align=center,
              text width=3.0cm, minimum height=1.5cm},
  ax/.style={-{Stealth[length=2.4mm]}, neutral, line width=0.9pt},
  tick/.style={font=\scriptsize, text=neutral, align=center}
]
  \node[opt=greenhl]   (a) at (0,0)    {Marketplace\\only\\{\scriptsize\mdseries
       somebody else's shop}};
  \node[opt=secondary] (b) at (3.7,0)  {Hosted\\platform\\{\scriptsize\mdseries
       rented, monthly}};
  \node[opt=primary]   (c) at (7.4,0)  {Self-hosted\\open source\\{\scriptsize
       \mdseries installed and extended}};
  \node[opt=purplehl]  (d) at (11.1,0) {Bespoke\\build\\{\scriptsize\mdseries
       written from nothing}};

  \draw[ax] (-1.8,-1.5) -- (12.9,-1.5);
  \node[tick] at (0,-1.95)    {no control\\no effort};
  \node[tick] at (11.1,-1.95) {total control\\total effort};
  \node[font=\scriptsize\itshape, text=accent] at (5.6,-2.55)
       {every step right costs more of your time and buys more of your fate};

  \node[font=\scriptsize\bfseries, text=primary] at (7.4,1.25)
       {this course};
  \draw[-{Stealth[length=1.8mm]}, primary, line width=0.9pt]
       (7.4,1.05) -- (7.4,0.80);
\end{tikzpicture}}%
{The four ways to obtain a shop. They are a spectrum, not a menu: each step to
the right buys control and costs effort, and the right answer depends on
which of the two you are shorter of.}{buildbuy}

\section{The Four Options, Honestly}

\rowcolors{2}{white}{lightbg}
\begin{longtable}{@{}L{2.8cm}L{5.8cm}L{6.0cm}@{}}
\toprule
\rowcolor{primary!15}
\textbf{Option} & \textbf{What it is good at} & \textbf{What it costs you} \\
\midrule
\endhead
\textbf{Marketplace only} & Customers on day one. No software, no hosting, no
payment integration. Genuinely the right answer for many small sellers & You
perform none of the four functions (\chref{what}). You do not own the customer
relationship, cannot email them, and can be delisted without appeal. The
commission is the visible cost and the smallest one \\
\textbf{Hosted platform} & Running in an afternoon, someone else patches it,
payment and shipping already integrated & A monthly fee that scales with your
success, a transaction cut on top, and a hard ceiling: when the platform
cannot do something, it cannot do it. Getting your data out ranges from
awkward to impossible \\
\textbf{Self-hosted open source} & You can change anything, including the
things the vendor never anticipated. No per-sale cut. The data is in your
database & You are now responsible for hosting, backups, updates and
\emph{security} --- an unpatched shop is a liability, not an inconvenience.
Needs someone who can do \chref{deploy} \\
\textbf{Bespoke build} & Fits the business exactly, because it was built for
it. The only option when the business is genuinely unusual & Twelve to
eighteen months before you sell anything, and you will reimplement baskets,
tax tables and payment retries that have been solved for twenty years ---
badly, because you will meet each edge case for the first time in production \\
\bottomrule
\end{longtable}

\begin{alertbox}[The outline this course follows assumes the fourth option]
The HEC outline for this course is, clause for clause, the contents page of a
2010 book about writing a shopping cart from scratch in PHP~5 --- see the
Course Specification. Taken literally, it asks a 5th-semester class to choose
\emph{bespoke build}.

\smallskip
That was defensible in 2010 and is not now, for a reason this chapter can
state precisely: the work has moved. In 2010, writing a basket and a checkout
was most of the job. Today a mature, free, auditable implementation of both is
one command away, and the job that remains --- modelling the catalogue,
configuring tax and shipping, integrating payment, securing the result,
extending what the platform cannot do --- is both larger and more
representative of what an IT graduate is actually hired to do.

\smallskip
So this course teaches every \emph{topic} the outline lists, and reaches them
by the third route rather than the fourth. You should be able to defend that
choice, not merely inherit it, which is what this chapter is for --- and
\chref{modules} is where the ``extend'' in ``self-hosted and extend'' stops
being a word and becomes code you write.
\end{alertbox}

\section{Requirements That Actually Discriminate}

Students write requirements like ``the system shall allow customers to add
products to a basket''. Every one of the four options does that. A requirement
that all four satisfy carries no information and cannot help you choose.

\begin{definitionbox}[Functional and Non-Functional Requirements]
A \term{functional requirement} says what the system must \emph{do}: accept a
voucher code, calculate tax by province, email an invoice.

\smallskip
A \term{non-functional requirement} says how well, under what constraints: how
fast, how many at once, how available, how secure, how maintainable, by whom.

\smallskip
For the build-buy-extend decision the useful requirements are nearly all
non-functional or unusual-functional. The ordinary functional ones are table
stakes --- everything does them --- and listing them at length is the most
common way a requirements document ends up being thirty pages that decide
nothing.
\end{definitionbox}

\begin{conceptbox}[The five questions that decide it]
In practice the choice turns on a handful of questions. Ask these first, and
most of the rest is detail.

\begin{tightnum}
  \item \textbf{Is there anything genuinely unusual about how you price, sell
        or deliver?} Per-customer negotiated pricing, made-to-measure items,
        rental with a return date, services booked by time slot. If yes, every
        option except the last two gets harder fast.
  \item \textbf{Who will maintain it in two years?} Not who will build it. A
        self-hosted shop with nobody to patch it becomes a breach
        (\chref{security}). If the honest answer is nobody, choose a hosted
        platform and accept the ceiling.
  \item \textbf{What volume, and how uneven?} Thirty orders a day and three
        thousand are different systems; so are steady traffic and traffic that
        multiplies tenfold for one week a year.
  \item \textbf{How much does leaving cost?} Can you export your products,
        customers and order history in a usable form? Ask it on day one,
        because the answer never improves.
  \item \textbf{What is the budget, over three years, including somebody's
        time?} Not the setup cost. See below.
\end{tightnum}
\end{conceptbox}

\section{The Cost That Is Not on the Price List}

The commonest error here is comparing a monthly fee against ``free'' and
concluding that open source is cheaper. Both numbers are wrong.

\begin{worked}[Three years of a small shop, two ways]
A shop expecting \textbf{Rs.~800{,}000 of sales a month}, growing slowly.
Compare a hosted platform against self-hosted open source over three years.

\smallskip
\textbf{Hosted platform.}

\begin{tight}
  \item Subscription: \$79/month $\approx$ Rs.~22{,}000. Over 36 months:
        \textbf{Rs.~792{,}000}.
  \item Platform transaction fee, 0.5\% of sales --- charged on top of the
        gateway's own fee: $0.005 \times 800{,}000 \times 36 =
        \textbf{Rs.~144{,}000}$.
  \item Two paid add-ons at \$15/month each $\approx$ Rs.~8{,}400/month:
        \textbf{Rs.~302{,}400}.
  \item Theme and setup, one-off: \textbf{Rs.~60{,}000}.
  \item Maintenance labour: near zero; the vendor patches.
  \item \textbf{Total: Rs.~1{,}298{,}400.}
\end{tight}

\smallskip
\textbf{Self-hosted open source.}

\begin{tight}
  \item Licence: \textbf{Rs.~0}. This is the number people stop at.
  \item Hosting, a modest virtual server plus backups, Rs.~9{,}000/month:
        \textbf{Rs.~324{,}000}.
  \item Setup and theming, 70 hours of somebody at Rs.~2{,}500/hour:
        \textbf{Rs.~175{,}000}.
  \item Two custom modules the platform does not ship, 40 hours:
        \textbf{Rs.~100{,}000}.
  \item \textbf{Maintenance: 4 hours a month} --- updates, patches, backups
        checked, the occasional breakage. $4 \times 36 \times 2{,}500 =
        \textbf{Rs.~360{,}000}$.
  \item \textbf{Total: Rs.~959{,}000.}
\end{tight}

\smallskip
\textbf{The result.} Self-hosting wins by Rs.~339{,}400 over three years ---
about 26\%. But look at what decided it: \textbf{maintenance labour is
Rs.~360{,}000}, the single largest line in the self-hosted column and the one
most often entered as zero. Set it to zero and self-hosting appears to win by
87\%. Double it, to eight hours a month because nobody on staff really knows
the platform, and the two options are level.

\smallskip
\textbf{The lesson.} This calculation is not really about software prices. It
is a question about whether you have, and will keep, somebody who can maintain
a server --- and if the answer is no, the hosted option is cheaper at any
subscription price, because the alternative is not a cheaper shop but an
unpatched one.
\end{worked}

\begin{alertbox}[Two costs that appear in neither column]
\textbf{Vendor lock-in} is not a line item until the day you want to leave, at
which point it is the whole bill. Before committing, actually export your
catalogue and customers from a trial account and look at what comes out. ``We
support CSV export'' and ``you can reconstruct your shop from this'' are
different claims.

\smallskip
\textbf{The ceiling} has no price at all. It is the feature you cannot have,
discovered in month fourteen, when the business depends on it. The hosted
option's real cost is the probability that this happens multiplied by what it
does to you --- which is why question~1 above is first.
\end{alertbox}

\section{Making the Decision Defensible}

Write it down as a matrix: criteria in rows, weights, options scored. Not
because the arithmetic is wise --- it is easy to fit the weights to the answer
you wanted --- but because it forces the weights into the open where somebody
can disagree with them.

\begin{examplebox}[A decision matrix for the shop in this course]
Criteria weighted out of 100, each option scored 1--5.

\rowcolors{2}{white}{lightbg}
\begin{longtable}{@{}L{4.8cm}L{1.5cm}L{1.9cm}L{1.9cm}L{1.9cm}@{}}
\toprule
\rowcolor{primary!15}
\textbf{Criterion} & \textbf{Weight} & \textbf{Hosted} & \textbf{Self-host} &
\textbf{Bespoke} \\
\midrule
\endhead
Students must see and change real code & 30 & 1 & 5 & 5 \\
Working shop within three weeks & 25 & 5 & 4 & 1 \\
Zero licence or subscription cost & 15 & 2 & 5 & 5 \\
Runs identically on every laptop & 15 & 3 & 5 & 4 \\
Covers the HEC clause list directly & 15 & 3 & 5 & 4 \\
\midrule
\textbf{Weighted total} & \textbf{100} & \textbf{2.75} & \textbf{4.75} &
\textbf{3.70} \\
\bottomrule
\end{longtable}

\smallskip
\textbf{Read the weights, not the totals.} The first row is worth 30 and it is
where hosted loses the contest --- a course whose students cannot read the code
is not this course. Change that weight to 5, as a small shopkeeper would, and
hosted wins outright. The matrix did not make the decision; it made the
decision \emph{visible}, which is all a matrix is for.
\end{examplebox}

\begin{pitfall}
\begin{tight}
  \item \textbf{Requirements that every option satisfies.} ``Customers can add
        to a basket'' discriminates between nothing.
  \item \textbf{Entering maintenance as zero.} It was the largest single line
        in the worked example, and leaving it out reverses the answer.
  \item \textbf{Comparing setup cost instead of three-year cost.} Subscriptions
        and labour both accumulate; neither appears in a setup quote.
  \item \textbf{Choosing bespoke because the platform is ``bloated''.} You will
        reimplement the parts you called bloat, in production, one incident at
        a time.
  \item \textbf{Choosing self-hosted with nobody to maintain it.} This is not a
        cheaper shop. It is a shop that will be compromised, and
        \chref{security} shows how quickly.
  \item \textbf{Not testing the export before committing.} Lock-in costs
        nothing until it costs everything.
  \item \textbf{Weighting the matrix after seeing the totals.} Everyone is
        tempted; write the weights down first and date them.
\end{tight}
\end{pitfall}

\begin{lab}[Your requirements, and your decision]
The last analysis laboratory. \chref{stack} installs software.

\smallskip
Using the venture from \chref{models}'s laboratory, produce a two-page
document:

\begin{tightnum}
  \item \textbf{Twelve requirements}, each labelled functional or
        non-functional, each numbered so you can refer to it later. At least
        four must be non-functional, with a number attached: not ``must be
        fast'' but ``a category page must render in under one second with
        5{,}000 products''. You will test some of these in
        \chref{catalogue} and \chref{deploy}, so make them checkable.
  \item \textbf{Mark the discriminating ones.} Go through your twelve and strike
        out every requirement that all four options satisfy. Most students
        strike out eight or nine. What is left is your real requirements
        document.
  \item \textbf{Answer the five questions} from this chapter, in writing,
        including the honest answer to ``who maintains this in two years?''
  \item \textbf{A three-year total cost of ownership}, two columns, following
        the worked example. Mark every figure you guessed, and state the one
        number the answer is most sensitive to.
  \item \textbf{A decision matrix}, weights written before scores, and a
        two-sentence conclusion.
  \item \textbf{The strongest objection to your decision}, and your answer to
        it. If you cannot state a real objection, you have not understood the
        option you rejected.
\end{tightnum}

\smallskip
\textbf{What to notice.} Your decision does not have to be the one this course
makes. A student whose venture is three products sold through Instagram should
conclude ``marketplace only'' and say so --- and will still build the shop in
Parts II--V, because the course has to teach the general case. Being right
about your own venture and building something else deliberately is a perfectly
coherent position, and a more useful one than pretending.
\end{lab}

\begin{checkpoint}
\begin{tightnum}
  \item Classify each as functional or non-functional, and say which of them
        discriminates between the four options: (a)~customers can reset a
        forgotten password; (b)~prices are agreed per customer and may differ
        for the same product; (c)~the shop stays up during a sale that brings
        twenty times normal traffic.
  \item A shop compares Rs.~22{,}000/month hosted against ``free'' open source
        and concludes open source saves Rs.~792{,}000 over three years. Name
        the three largest costs they have omitted.
  \item Give a specific business for which this course's choice --- self-hosted
        open source --- would be the wrong answer, and say which option you
        would choose instead and why.
\end{tightnum}
\end{checkpoint}

\begin{chaptersummary}
\begin{tight}
  \item A shop can be obtained four ways: marketplace only, hosted platform,
        self-hosted open source, or bespoke build. They form a spectrum trading
        control against effort.
  \item Requirements that every option satisfies cannot help you choose. Strike
        them out; what remains is the document that matters.
  \item Five questions decide most cases: anything unusual about pricing or
        delivery, who maintains it in two years, volume and its unevenness, the
        cost of leaving, and three-year budget including labour.
  \item Compare total cost of ownership over three years, not setup cost. In
        the worked example maintenance labour was the largest self-hosted line,
        and entering it as zero reverses the conclusion.
  \item Vendor lock-in and the feature ceiling appear in no price list and are
        often the real cost of the hosted option.
  \item A decision matrix does not make the decision; it makes the weights
        visible so somebody can argue with them. Write the weights first.
  \item This course chooses self-hosted open source and teaches every clause of
        the outline through it, rather than writing a cart from scratch as the
        2010 book behind the outline assumed. You should be able to defend
        that, including knowing when it is the wrong answer.
\end{tight}
\end{chaptersummary}

\begin{reviewq}
  \item \textbf{(MCQ)} Which is a non-functional requirement? (a)~The shop
        emails an invoice after each order; (b)~the checkout completes in under
        two seconds at 200 concurrent users; (c)~customers can apply a voucher
        code; (d)~administrators can issue refunds.
  \item \textbf{(MCQ)} In the worked example, the largest single cost in the
        self-hosted column was: (a)~hosting; (b)~setup and theming; (c)~custom
        modules; (d)~maintenance labour.
  \item \textbf{(Short)} Explain why a requirement satisfied by all four options
        is useless for this decision, and give an example of one.
  \item \textbf{(Short)} What is vendor lock-in, and what concrete test would
        you run during a free trial to measure it?
  \item \textbf{(Short)} Give the strongest argument \emph{against} the choice
        this course makes, and then answer it.
  \item \textbf{(Applied)} A tailoring business takes measurements in the shop,
        makes garments to order over two weeks, and prices each job
        individually after seeing the cloth. Work through the five questions
        and recommend one of the four options, justifying it. Identify the one
        requirement that makes this case different from a normal shop.
  \item \textbf{(Applied)} Build a three-year total cost of ownership for a
        shop doing Rs.~3{,}000{,}000 a month --- nearly four times the worked
        example --- on both options, keeping all other assumptions. State which
        option wins, and identify the single assumption that, if changed, would
        reverse it.
\end{reviewq}
